\scanpage{86}
% archive foliation: 88
qui opère sur $\operatorname{Oct}(V)$, opère sur $I(V)$, et on trouve
\diagram{\begin{tikzcd}[ampersand replacement=\&]
\Gamma(V) \arrow[r,"\sim"] \arrow[d,equal] \& \unsure{\mathfrak S}_{I(V)},\\
\operatorname{Aut}^+(\operatorname{Oct}(V)) \&
\end{tikzcd}}
et cet iso \deleted{$\unsure{\operatorname{Aut}}$} définit une structure de
$V$-torseur sur $I(V)$. (Ceci résulte de la théorie
combinatoire de l'octaèdre, ou du cube c'est kif kif.)

\uline{Proposition} Soit $X$ droite projective sur un corps alg. clos $k$
(il suffit qu'il \deleted{contienne} soit stable par extraction des racines
carrées) de car. $\ne2$, et $G=\operatorname{Aut}_kX$. Alors il y a corr 1-1
entre

a) les sous-groupes $V$ de $G$ isom. à $D_2$

b) \rule{3cm}{0.2pt} $\Gamma$ de $G$ isom. à $\unsure{\mathfrak S}_4$ \deleted{\unsure{un}}

c) les parties $S\in\unsure{P|\mathcal P}_6(X(k))$ qui sont ``octaédrales'', i.e.
\deleted{\unsure{telles que}}\inserted{isom. à} $(\mathbb P^1_k,\{0,\infty,\pm1,\pm i\})$

par
\[
\begin{aligned}
\Gamma&=\operatorname{Nor}_G(V)\\
V&=\inserted{\text{unique}}\ \text{sous-groupe }\deleted{\text{\unsure{invariant isom. à}}}\inserted{\text{inv. de }\Gamma\ \text{isom. à}}\ D_2\ \deleted{\unsure{\subset G}}\\
S&=\text{ens de ramification de l'action de }V\text{ sur }X
\end{aligned}
\]
\deleted{\unsure{ou de l'action}}
(\uline{NB} $\Gamma$ a un ens de ramification
plus grand qui comprend, en plus de $S$,
les ``milieux d'arêtes'' et ``centres des faces''
de l'octaèdre,

\margin{\uline{Cor} la catégorie (quand $k$ fixé)
des couples $(X,V)$ (ou $(X,\Gamma)$
$(X,S)$) $\unsure{\longrightarrow|\text{est}}$ équivalente à celle
des octaèdres \inserted{\unsure{combin.}} \uline{orientés}, une
fois choisi \deleted{\unsure{a}} une racine
carrée de $-1$ dans $k$. Mieux : la
cat. des couples $(X,V)$ $\unsure{\longrightarrow|\text{est}}$ équivalente
à celle des couples $(\operatorname{Oc},\varphi)$
où $\operatorname{Oc}$ est un octaèdre combinatoire et
$\varphi:\{\pm i\}\overset\sim\longrightarrow\unsure{\operatorname{or}}(\operatorname{Oc})$.}

Dans l'ensemble $S$, on a un antipodisme $\uline{a}_S=\uline{a}_V$
défini par la condition que pour $\sigma\in V^*$, $a_S$ \deleted{est}\inserted{induit} la transposition
de $X^\sigma$ (de cardinal 2). On définit sur $S$ une
structure octaédrale, admettant $\Gamma$ comme groupe des
autom. orientés, et $\uline a$ comme antipodisme, via l'antipodisme
(les faces sont les parties $S'\in\unsure{P|\mathcal P}_3(S)$ telles que $S'\cup-S'=\unsure{S'|S}$).
